% Generated draft: requires mathematical and source review.
\documentclass[11pt]{article}
\usepackage[margin=25mm]{geometry}
\usepackage{fontspec}
\setmainfont{Latin Modern Roman}
\usepackage{amsmath,amssymb,mathtools,unicode-math}
\setmathfont{Latin Modern Math}
\usepackage{enumitem,needspace,xurl,hyperref}
\setlength{\parindent}{0pt}
\setlength{\parskip}{5pt}
\setlength{\emergencystretch}{3em}
\newcommand{\N}{\mathbb{N}}
\newcommand{\Z}{\mathbb{Z}}
\newcommand{\Q}{\mathbb{Q}}
\newcommand{\R}{\mathbb{R}}
\newcommand{\C}{\mathbb{C}}
\newcommand{\F}{\mathbb{F}}
\newcommand{\A}{\mathbb{A}}
\newcommand{\PP}{\mathbb{P}}
\newcommand{\E}{\mathbb{E}}
\newcommand{\eps}{\varepsilon}
\newcommand{\dd}{\,\mathrm{d}}
\begin{document}
\section*{UnsolvedMath ID: 2308019}
\textbf{Source problem code:} catalogue entry
\subsection{Definitions and mathematical statement}
For a finite positive Borel measure $\mu$ on $\mathbb T$ singular with respect to arclength, the supplied normalized singular inner function is
\[
 S_\mu(z)=\exp\left\{\int_0^{2\pi}
       \frac{z+e^{i\theta}}{z-e^{i\theta}}\,\dd\mu(\theta)\right\},
 \qquad z\in\mathbb D.
\]
In particular $S_\mu(0)=e^{-\mu(\mathbb T)}$. It is called discrete if $\mu$ is purely atomic, and continuous if $\mu$ has no atoms. Here continuous refers to the measure type; it does not mean that the inner function extends continuously to the closed disk.

K. Stephenson's question asks whether every admissible fixed discrete $S_\mu$ is separated in the $H^\infty$ norm from all continuous singular inner functions:
\[
 \exists\delta=\delta(\mu)>0\ \forall\nu\text{ admissible atomless singular measures},
 \qquad\|S_\mu-S_\nu\|_\infty>\delta,
\]
where $\|g\|_\infty=\sup_{z\in\mathbb D}|g(z)|$. The quantifier is for one fixed $\mu$ at a time; a uniform lower bound for all discrete measures is not asserted. Only the normalized functions in the displayed construction are used, without freely multiplying by unimodular constants.

The source begins ``$\mu\geq1$ be a singular measure,'' which does not specify whether a total-mass normalization is intended and is not a standard unqualified ordering of a measure against the number one. Finite positivity is the usual requirement for the displayed singular-inner construction, but it does not resolve that extra source condition. Exclusion of the zero measure and the allowed total masses for $\nu$ also need confirmation; the zero measure would yield the constant function one and belongs to both measure types if permitted.

The source reports a negative answer when discrete and continuous are interchanged. This is retained as failure of the corresponding universal separation statement, not strengthened into an assertion that every atomless singular inner function admits a particular atomic approximation. The lowercase $s_\nu$ in the printed norm is interpreted as its $S_\nu$; no separate lowercase family is defined.
\subsection{Short English statement}
Is each discrete singular inner function a positive uniform distance from the family arising from atomless singular measures? The measure normalization in the source is unclear, and continuous describes atomlessness rather than boundary continuity. The source reports that the reversed separation question has a negative answer. Its supplied triage says the original was open in Hayman's 2018 update and no definitive resolution was located; no independent current-status claim is made.
\subsection{Sources}
\begin{itemize}
\item UnsolvedMath permanent ID 2308019, source code catalogue entry, ``Research Problems in Function Theory --- Problem 8.19.'' Registry curation and metadata: Ulam AI and UnsolvedMath Contributors, CC BY 4.0; updated 2026-08-06.
\item Hayman, \emph{Research Problems in Function Theory} (2018 source update), Problem 8.19, attributed to K. Stephenson: \url{https://arxiv.org/abs/1809.07200}. Supplied source-TeX extraction accessed 2026-07-29. The metadata retains the update for status triage and marks status/release and rights review NEEDS\_REVIEW.
\item No complete separate reference is supplied for the negative reversed statement. No missing author, title, example or locator is inferred.
\end{itemize}
\end{document}
